\chapter{Templates : écrire une fois pour tous les types}

Un \emph{template} --- un \textbf{modèle} --- écrit du code une seule fois, et le
compilateur le décline pour chaque type qu'on lui demande. C'est ce qui permet à
\texttt{NkVector} d'exister sans être réécrit pour les entiers, puis pour les
chaînes, puis pour les entités.

Ce chapitre va du modèle de fonction le plus simple jusqu'aux traits de type ---
la mécanique qui laisse \texttt{NkVector} choisir entre copier et déplacer selon
ce qu'il transporte. C'est la partie du C++ qui fait peur au début ; elle repose
en réalité sur trois idées seulement.

\section{Modèle de fonction}

\begin{nkcode}{Exemple pedagogique}
template <typename T>
T Maximum(const T &a, const T &b) {
    return (a > b) ? a : b;
}

int32   x = Maximum(3, 7);        // T = int32, DEDUIT des arguments
float32 y = Maximum(1.f, 2.f);    // T = float32
float32 z = Maximum<float32>(1.f, 2);   // T IMPOSE : le 2 sera converti
\end{nkcode}

Le compilateur \textbf{déduit} \texttt{T} des arguments. Quand il ne peut pas, ou
quand la déduction donnerait le mauvais type, on l'impose entre chevrons.

\begin{nkretenir}
Un template n'est pas du code : c'est une \textbf{recette}. Rien n'est compilé
tant que personne ne l'emploie. \texttt{Maximum} écrit et jamais appelé ne
produit pas une seule instruction --- et ses erreurs ne sont pas détectées non
plus.
\end{nkretenir}

\begin{nkpiege}
\textbf{La déduction ne convertit pas.} \texttt{Maximum(1.f, 2)} échoue : le
premier argument dit \texttt{T = float32}, le second dit \texttt{T = int32}, et
le compilateur refuse d'arbitrer. Il faut imposer \texttt{Maximum<float32>(1.f,
2)}, ou corriger l'appel.

Le message d'erreur parle de « deduced conflicting types » : c'est lui, et il
est plus clair qu'il n'en a l'air.
\end{nkpiege}

\section{Modèle de classe}

\begin{nkcode}{Principe de NkVector, simplifie}
template <typename T>
class NkTableau {
    public:
        void PushBack(const T &v);
        T &operator[](usize i);
        usize Size() const;

    private:
        T *mDonnees = nullptr;
        usize mTaille = 0;
        usize mCapacite = 0;
};

NkTableau<int32>  entiers;
NkTableau<NkVec2f> points;   // une AUTRE classe, engendree du meme modele
\end{nkcode}

\texttt{NkTableau<int32>} et \texttt{NkTableau<NkVec2f>} sont deux classes
\emph{distinctes}, sans rapport de parenté. Elles viennent de la même recette,
c'est tout.

\begin{nkpiege}
\textbf{Le code d'un template vit dans l'en-tête}, jamais dans un \texttt{.cpp}.
Le compilateur doit \emph{voir} le corps pour l'instancier avec votre type ; s'il
est dans un \texttt{.cpp} qu'il ne lit pas, il croit la déclaration et laisse le
corps au lieur --- qui ne le trouve nulle part.

Symptôme : \texttt{undefined reference to NkTableau<int>::PushBack}. Le code
compile, il ne \emph{lie} pas. C'est pourquoi \texttt{NkVector.h} fait
2\,500 lignes et n'a pas de \texttt{.cpp}.
\end{nkpiege}

\section{Paramètres qui ne sont pas des types}

Un modèle peut prendre une \emph{valeur}, connue à la compilation.

\begin{nkcode}{Exemple pedagogique}
template <typename T, usize N>
class NkTableauFixe {
    public:
        usize Size() const { return N; }
    private:
        T mDonnees[N];        // taille connue A LA COMPILATION : aucune allocation
};

NkTableauFixe<float32, 16> matrice;
\end{nkcode}

\begin{nkretenir}
C'est ainsi qu'un tableau de taille fixe garde sa taille --- contrairement au
tableau C du chapitre précédent, qui la perdait en passant à une fonction. La
taille fait partie du \emph{type}.
\end{nkretenir}

\section{Spécialisation}

Parfois un type précis demande un traitement à part. On \emph{spécialise}.

\begin{nkcode}{Exemple pedagogique}
template <typename T>
bool EstVide(const T &v) { return v.Size() == 0; }

// Specialisation TOTALE : ce cas-la, et lui seul, se traite autrement.
template <>
bool EstVide<const char *>(const char *const &s) { return s == nullptr || *s == 0; }
\end{nkcode}

\begin{nkpiege}
\textbf{Une spécialisation qui ne correspond pas exactement est ignorée en
silence} --- le modèle général est employé à sa place, et votre cas particulier
n'est jamais pris. Le programme compile et se comporte mal.

Vérifiez la signature au caractère près : \texttt{const char *} et
\texttt{char *} sont deux types différents.
\end{nkpiege}

\section{Modèles variadiques}

Un modèle peut prendre un nombre \emph{quelconque} d'arguments.

\begin{nkcode}{Principe de EmplaceBack}
template <typename... Args>
void EmplaceBack(Args &&...args) {
    // construit l'element DIRECTEMENT dans le tableau,
    // au lieu de le construire puis de le copier
    new (Emplacement()) T(NkForward<Args>(args)...);
}
\end{nkcode}

\texttt{Args...} est un \emph{paquet} de types ; \texttt{args...} le paquet de
valeurs. Les trois points « déplient » le paquet à l'endroit où on les écrit.

\begin{nkretenir}
C'est la différence entre \texttt{PushBack} et \texttt{EmplaceBack} :
\texttt{PushBack} reçoit un objet \emph{déjà construit} et le copie ou le
déplace ; \texttt{EmplaceBack} reçoit les \emph{arguments} et construit l'objet
sur place. Une construction au lieu de deux.
\end{nkretenir}

\section{Références universelles et transfert parfait}

\begin{nkcode}{Exemple pedagogique}
template <typename T>
void Relayer(T &&arg) {           // `T&&` dans un TEMPLATE : reference universelle
    Consommer(NkForward<T>(arg)); // conserve : lvalue reste lvalue, rvalue reste rvalue
}
\end{nkcode}

\begin{nkpiege}
\textbf{\texttt{T\&\&} ne veut pas dire la même chose selon le contexte.} Dans un
template avec \texttt{T} déduit, c'est une \emph{référence universelle} : elle
accepte les deux. Hors template --- \texttt{NkVector\&\&} --- c'est une vraie
référence rvalue, qui n'accepte que les temporaires.

Deux notions, une seule écriture. C'est le point du langage qui coûte le plus de
malentendus.
\end{nkpiege}

\begin{nkpiege}
\textbf{\texttt{NkMove} et \texttt{NkForward} ne déplacent rien.} Ce sont des
\emph{conversions de type} : elles disent au compilateur « traite ceci comme un
temporaire ». Le déplacement, s'il a lieu, est fait par le constructeur de
déplacement de la classe.

Les compter comme des allocations ou des corruptions mémoire est faux. Ce qu'ils
signalent, c'est du code qui raisonne en durée de vie --- une information, pas un
défaut.
\end{nkpiege}

\section{Traits de type}

Un \emph{trait} répond à une question sur un type, à la compilation.

\begin{nkcode}{Principe employe par NkVector}
if (traits::NkIsTriviallyCopyable<T>::value) {
    // un memcpy suffit : le type n'a pas de constructeur a respecter
} else {
    // il faut construire chaque element un par un
}
\end{nkcode}

C'est ainsi que \texttt{NkVector} peut recopier mille \texttt{int32} d'un bloc,
et mille \texttt{NkString} un par un --- sans que vous ayez à le lui dire.

\begin{nkretenir}
\textbf{Tout se décide à la compilation.} Le test ci-dessus ne coûte rien à
l'exécution : le compilateur sait déjà quelle branche garder, et supprime
l'autre. Un trait n'est pas un \texttt{if} déguisé, c'est un choix fait avant que
le programme existe.
\end{nkretenir}

\section{Ce que les templates coûtent}

\begin{center}
\begin{tabular}{@{}ll@{}}
\toprule
\textbf{Avantage} & \textbf{Prix} \\
\midrule
un seul code pour tous les types & le binaire grossit : une copie par type \\
aucun coût à l'exécution & la compilation ralentit \\
erreurs vues à la compilation & messages longs et indirects \\
\bottomrule
\end{tabular}
\end{center}

\begin{nkpiege}
\textbf{Lisez la PREMIÈRE ligne d'une erreur de template.} Elle nomme la vraie
cause. Les vingt suivantes déroulent la pile d'instanciation : elles disent
\emph{par où} le compilateur est arrivé là, pas ce qui ne va pas.

Le réflexe qui fait gagner le plus de temps : chercher, dans le pavé, la
première ligne qui cite \emph{votre} fichier.
\end{nkpiege}

\begin{nkbilan}
Un template est une \textbf{recette}, pas du code : rien n'existe avant qu'on
l'emploie. Son corps vit dans l'en-tête, sinon l'édition de liens échoue. Il se
décline par type (\texttt{typename T}) ou par valeur (\texttt{usize N}), se
spécialise pour un cas particulier, et se généralise à un nombre quelconque
d'arguments. Les traits lui permettent de choisir sa stratégie à la compilation
--- ce qui rend \texttt{NkVector} rapide sur les types simples sans être faux sur
les autres.
\end{nkbilan}

\begin{nkquestions}
\begin{enumerate}
\item Pourquoi le corps d'un template doit-il vivre dans l'en-tête ?
\item Quelle différence entre \texttt{PushBack} et \texttt{EmplaceBack} ?
\item \texttt{T\&\&} dans un template et \texttt{NkVector\&\&} hors template :
      qu'est-ce qui les distingue ?
\item \texttt{NkMove} déplace-t-il quelque chose ? Justifiez.
\item Un template jamais employé peut-il contenir une erreur ? Pourquoi ?
\end{enumerate}
\end{nkquestions}

\begin{nkpratique}
Écrivez \texttt{NkPaire<A, B>} : deux champs publics, un constructeur, et une
méthode \texttt{Echanger()} qui \emph{rend} une \texttt{NkPaire<B, A>}.

Puis ajoutez un modèle de fonction \texttt{Fabriquer(a, b)} qui construit la
paire en déduisant les deux types --- de sorte qu'on écrive
\texttt{Fabriquer(3, 1.f)} sans chevrons.
\end{nkpratique}

\begin{nkdemo}
Prouvez, sans lire la documentation, que le corps d'un template doit être dans
l'en-tête.

Écrivez \texttt{NkBoite<T>} avec une méthode déclarée dans le \texttt{.h} et
\emph{définie dans un \texttt{.cpp}}. Employez-la depuis un autre fichier.
Recopiez le message d'erreur exact, et dites à quelle étape il survient ---
compilation ou édition de liens. Déplacez ensuite la définition dans l'en-tête et
constatez.
\end{nkdemo}

\begin{nklogique}
\texttt{NkVector<int32>} et \texttt{NkVector<float32>} sont deux classes
distinctes engendrées du même modèle.

Une fonction \texttt{void Afficher(NkVector<int32> \&v)} peut-elle recevoir un
\texttt{NkVector<float32>} ? Et si l'on écrit
\texttt{template <typename T> void Afficher(NkVector<T> \&v)} ?

Déduisez-en ce que le compilateur engendre réellement, et pourquoi un binaire qui
emploie dix types différents dans un \texttt{NkVector} contient dix jeux de code.
\end{nklogique}
